\let\negmedspace\undefined
\let\negthickspace\undefined
\documentclass[journal,onecolumn]{IEEEtran}

\usepackage[a5paper, margin=10mm]{geometry} % Safe with onecolumn
\usepackage{lmodern} % Better fonts (optional but recommended)
\usepackage{tfrupee} % Rupee symbol package

% Header fix
\setlength{\headheight}{12pt} % Safe value
\setlength{\headsep}{0pt}     % Reduce gap

\usepackage{gvv-book}
\usepackage{gvv}
\usepackage{cite}
\usepackage{amsmath,amssymb,amsfonts,amsthm}
\usepackage{algorithmic}
\usepackage{graphicx}
\usepackage{textcomp}
\usepackage{xcolor}
\usepackage{txfonts}
\usepackage{listings}
\usepackage{enumitem}
\usepackage{mathtools}
\usepackage{gensymb}
\usepackage{comment}
\usepackage[breaklinks=true]{hyperref}
\usepackage{tkz-euclide} 
\usepackage{listings}
% \usepackage{gvv}                                        
\def\inputGnumericTable{}                                 
\usepackage[latin1]{inputenc}                                
\usepackage{color}                                            
\usepackage{array}                                            
\usepackage{longtable}                                       
\usepackage{calc}                                             
\usepackage{multirow}                                         
\usepackage{hhline}                                           
\usepackage{ifthen}                                           
\usepackage{lscape}

\usepackage{amsmath,amssymb}
\usepackage{booktabs}
\usepackage{tikz}
\usetikzlibrary{arrows.meta,angles,quotes}

\begin{document}

\bibliographystyle{IEEEtran}
\vspace{3cm}

\title{5.4.9}
\author{EE25BTECH11014 - Bhoomika Lokesh}
% \maketitle
% \newpage
% \bigskip
{\let\newpage\relax\maketitle}

\renewcommand{\thefigure}{\theenumi}
\renewcommand{\thetable}{\theenumi}
\setlength{\intextsep}{10pt} % Space between text and floats

% Remove this line so equations use default numbering (1, 2, 3,...)
% \numberwithin{equation}{enumi}
\numberwithin{figure}{enumi}
\renewcommand{\thetable}{\theenumi}


\textbf{Question:}Using elementary transformations, find the inverse of the following matrix\\
\begin{center}
    $\myvec{1&2\\3&4}$
\end{center}
\textbf{Solution:}\\
We know that,\\
\begin{align}
    \vec{A^{-1}}\vec{A}=\vec{I}
\end{align}
    where $\vec{I}$ is a 2 X 2 identity matrix.\\
    The following is the augmented matrix we get,
\begin{align}
  \myvec{
    \begin{array}{cc|cc}
      1 & 2 & 1 & 0 \\ 
      3 & 4 & 0 & 1
    \end{array}
  }
  \xrightarrow{\;R_2\longrightarrow R_2-3R_1\;}   \myvec{
    \begin{array}{cc|cc}
      1 & 2 & 1 & 0 \\ 
      0 & -2 & -3 & 1
    \end{array}
  }\\
  \myvec{
    \begin{array}{cc|cc}
      1 & 2 & 1 & 0 \\ 
      0 & -2 & -3 & 1
    \end{array}
  }\xrightarrow{\;R_2\longrightarrow \frac{R_2}{-2}\;} 
  \myvec{
    \begin{array}{cc|cc}
      1 & 2 & 1 & 0 \\ 
      0 & 1 & \frac{3}{2} & \frac{-1}{2}
    \end{array}
  }\\
   \myvec{\begin{array}{cc|cc}
      1 & 2 & 1 & 0 \\ 
      0 & 1 & \frac{3}{2} & \frac{-1}{2}
    \end{array}
  }\xrightarrow{\;R_1\longrightarrow R_1-2R_2\;}   \myvec{\begin{array}{cc|cc}
      1 & 0 & -2 & 1 \\ 
      0 & 1 & \frac{3}{2} & \frac{-1}{2}
 \end{array}}
\end{align}
\centering
    $\therefore \vec{A^{-1}}=\frac{1}{2}\myvec{-4 & 2\\3 & -1}$


\end{document}